\subsection{Groups of Permutations}

Before diving deeper into permutations, we introduce a fundamental concept that relates different groups by preserving their mathematical structure.


%%%%%%%%%%%%%%%%%%%
%  Homomorphisms  %
%%%%%%%%%%%%%%%%%%%

\subsubsection{Homomorphisms}

Recall that an isomorphism between two groups is a bijection that preserves the group operation. A homomorphism is a more general mapping that preserves the group operation but is not necessarily one-to-one or onto.

\Definition{Homomorphism}{
    Let $G$ and $G'$ be groups with a function $\phi : G \to G'$. The map $\phi$ is a \textbf{homomorphism} if the homomorphism property \[
        \phi(ab) = \phi(a)\phi(b)
    \] holds for all $a, b \in G$.

    Here, the left-hand side uses the group operation in $G$, while the right-hand side uses the group operation in $G'$. If $\phi$ is a homomorphism that is also a bijection, then $\phi$ is an isomorphism.
}

\Example{}{
    Let $\phi: \R \to U$ (the circle group) be defined by $\phi(x) = e^{2\pi ix}$. Since $e^{2\pi i(a+b)} = e^{2\pi ia}e^{2\pi ib}$, we have $\phi(a+b) = \phi(a)\phi(b)$. Thus, $\phi$ is a homomorphism. It maps $\R$ onto $U$, but is not one-to-one (e.g., $\phi(1) = \phi(0) = 1$), so it is not an isomorphism.
}

\Definition{Image and Inverse Image}{
    If $\phi: X \to Y$ is a function, $A \subseteq X$, and $B \subseteq Y$, we define the image of $A$ as $\phi[A] = \{\phi(a) \mid a \in A\}$ and the inverse image of $B$ as $\phi^{-1}[B] = \{a \in A \mid \phi(a) \in B\}$.
}

\Theorem{Properties of Homomorphisms}{
    Let $\phi$ be a homomorphism of a group $G$ into a group $G'$.
    \begin{enumerate}
        \item If $e$ is the identity in $G$, then $\phi(e)$ is the identity $e'$ in $G'$.
        \item If $a \in G$, then $\phi(a^{-1}) = \phi(a)^{-1}$.
        \item If $H$ is a subgroup of $G$, then $\phi[H]$ is a subgroup of $G'$.
        \item If $K'$ is a subgroup of $G'$, then $\phi^{-1}[K']$ is a subgroup of $G$.
    \end{enumerate}
    Loosely speaking, homomorphisms preserve the identity, inverses, and subgroups.
}

\Proof{
    \begin{enumerate}
        \item $\phi(e) = \phi(ee) = \phi(e)\phi(e)$, so $\phi(e)$ is the identity in $G'$.
        \item $\phi(a)\phi(a^{-1}) = \phi(aa^{-1}) = \phi(e) = e'$, so $\phi(a^{-1})$ is the inverse of $\phi(a)$.
        \item If $x, y \in \phi[H]$, then $x = \phi(h_1)$ and $y = \phi(h_2)$ for some $h_1, h_2 \in H$. Since $H$ is a subgroup, $h_1 h_2^{-1} \in H$. Then $xy^{-1} = \phi(h_1)\phi(h_2)^{-1} = \phi(h_1)\phi(h_2^{-1}) = \phi(h_1 h_2^{-1})$, which is in $\phi[H]$. Thus, $\phi[H]$ is a subgroup of $G'$.
        \item If $x, y \in \phi^{-1}[K']$, then $\phi(x), \phi(y) \in K'$. Since $K'$ is a subgroup, $\phi(x)\phi(y)^{-1} = \phi(x)\phi(y^{-1}) = \phi(xy^{-1})$ is in $K'$. Thus, $xy^{-1}$ is in $\phi^{-1}[K']$, and $\phi^{-1}[K']$ is a subgroup of $G$.
    \end{enumerate}
}

\Definition{Kernel}{
    Let $\phi : G \to G'$ be a homomorphism of groups. The subgroup $\phi^{-1}[\{e'\}] = \{x \in G \mid \phi(x) = e'\}$ is called the \textbf{kernel} of $\phi$, denoted by $\ker(\phi)$.
}


%%%%%%%%%%%%%%%%%%%%%%
%  Cayley's Theorem  %
%%%%%%%%%%%%%%%%%%%%%%

\subsubsection{Cayley's Theorem}

Every group is structurally related to a group of permutations. The connection relies on the group table: each row represents a permutation of the group elements because multiplying the entire group by a fixed element $x$ merely permutes the elements.

\Definition{Left Regular Representation}{
    Let $G$ be a group. The function $\phi : G \to S_G$ given by $\phi(x) = \lambda_x$, where $\lambda_x(g) = xg$ for all $g \in G$, is called the \textbf{left regular representation} of $G$.
}

\Theorem{Cayley's Theorem}{
    Every group is isomorphic to a group of permutations.
}

\Proof{
    Let $G$ be a group and let $\phi : G \to S_G$ be the left regular representation defined by $\phi(x) = \lambda_x$. We must show that $\phi$ is a one-to-one homomorphism; then $\phi[G]$ is a subgroup of $S_G$, making $G$ isomorphic to the permutation group $\phi[G]$.

    First, suppose $\phi(a) = \phi(b)$. Then $\lambda_a = \lambda_b$, so $\lambda_a(e) = \lambda_b(e)$. This means $ae = be$, so $a = b$. Thus, $\phi$ is one-to-one.

    Next, we show $\phi$ is a homomorphism: we want $\phi(ab) = \phi(a)\phi(b)$, or $\lambda_{ab} = \lambda_a \lambda_b$. For any $g \in G$: \[
        \lambda_{ab}(g) = (ab)g = a(bg) = \lambda_a(bg) = \lambda_a(\lambda_b(g)) = (\lambda_a \circ \lambda_b)(g)
    \]
    Thus $\phi$ is a homomorphism, and the theorem is proven.
}

\Definition{Right Regular Representation}{
    Let $G$ be a group. The map $\tau : G \to S_G$ given by $\tau(x) = \sigma_{x^{-1}}$ where $\sigma(x) = gx$ is called the \textbf{right regular representation} of $G$.
}


%%%%%%%%%%%%%%%%%%%%%%%%%%%%%%%
%  Even and Odd Permutations  %
%%%%%%%%%%%%%%%%%%%%%%%%%%%%%%%

\subsubsection{Even and Odd Permutations}

\Definition{Transposition}{
    A cycle of length 2 is called a \textbf{transposition}. A transposition leaves all elements except two fixed and swaps those two.
}

A simple computation shows that any cycle of length $n$ can be decomposed into transpositions: \[
    (a_1, a_2, \ldots, a_n) = (a_1, a_n)(a_1, a_{n-1}) \cdots (a_1, a_3)(a_1, a_2)
\]
Because every permutation of a finite set is a product of disjoint cycles, we establish:

\Theorem{}{
    Any permutation of a finite set containing at least two elements can be written as a product of transpositions.
}

\Note{
    Any permutation of length $n$ can be written as a product of at most $n-1$ transpositions, since a cycle of length $n$ can be decomposed into $n-1$ transpositions.
}

\Note{
    This decomposition is not unique. For instance, the identity can be written as $(1,2)(1,2)$ or $(1,3)(1,3)(2,4)(2,4)$. A given permutation can therefore be written as the product of different numbers of transpositions. However, the \textit{parity} of the number of transpositions is an invariant.
}

\Definition{Orbit of a Permutation}{
    Let $\sigma$ be a permutation of a set $A$. The \textbf{orbit} of an element $a \in A$ under $\sigma$ is the set $\{\sigma^n(a) \mid n \in \Z\}$.
}

\Example{}{
    Let $\sigma=(1,2,6)(3,5)\in S_6$. Then the orbit of $1$, $2$, and $6$ is $\{ 1,2,6 \}$, the orbit of $3$ is $\{ 3,5 \}$, and the orbit of $4$ is $\{ 4 \}$ since $4$ is fixed by $\sigma$.
}

\Theorem{}{
    No permutation in the symmetric group $S_n$ can be expressed both as a product of an even number of transpositions and as a product of an odd number of transpositions.
}

\Proof{
    Let $\sigma$ be a permutation of $S_n$. We will show that the parity of the number of transpositions in any decomposition of $\sigma$ is invariant.

    Consider the action of $\sigma$ on the set $\{1, 2, \ldots, n\}$. The number of orbits under this action is equal to the number of cycles in the cycle decomposition of $\sigma$. Each cycle of length $k$ can be written as a product of $k-1$ transpositions. Therefore, if $\sigma$ has $c$ cycles, it can be expressed as a product of $n - c$ transpositions.

    Since $n - c$ is fixed for a given $\sigma$, the parity of the number of transpositions in any decomposition of $\sigma$ is also fixed. Thus, if $\sigma$ can be written as a product of an even number of transpositions, it cannot be written as a product of an odd number of transpositions, and vice versa.
}

\Definition{Even and Odd Permutations}{
    A permutation is \textbf{even} if it can be written as a product of an even number of transpositions, and \textbf{odd} if it can be written as a product of an odd number of transpositions.
}

\Example{}{
    In $S_3$, the transpositions are $(1,2)$, $(1,3)$, and $(2,3)$, which are odd. The 3-cycles $(1,2,3)$ and $(1,3,2)$ can be written as products of two transpositions (e.g., $(1,2,3) = (1,3)(1,2)$), so they are even. The identity is also even since it can be written as a product of zero transpositions. Thus, $S_3$ has three odd permutations and three even permutations.
}


%%%%%%%%%%%%%%%%%%%%%%%%
%  Alternating Groups  %
%%%%%%%%%%%%%%%%%%%%%%%%

\subsubsection{Alternating Groups}

\Theorem{}{
    If $n \geq 2$, then the collection of all even permutations of $\{ 1,2,3, \dots, n \}$ forms a subgroup of order $n!/2$ of the symmetric group $S_n$.
}

\Proof{
    Let $A_n$ be the set of all even permutations in $S_n$. We will show that $A_n$ is a subgroup of $S_n$ and that its order is $n!/2$.

    First, $A_n$ is closed under the group operation since the product of two even permutations is even. Second, if $\sigma$ is an even permutation, then its inverse $\sigma^{-1}$ is also even, because the inverse of a product of transpositions is the product of the inverses of those transpositions, and the number of transpositions in the inverse is the same as in the original permutation. Thus, $A_n$ contains inverses. So, $A_n$ is a subgroup of $S_n$.

    Finally, we need to show that $A_n$ is exactly half of $S_n$. For that, we can define a function called the \textbf{sign of a permutation}, $\text{sgn}: S_n \to \{ 1,-1 \}$ such that \[ 
        \text{sgn}(\sigma) = \begin{cases}
            1, &\text{ if $\sigma$ is even}\\
            -1, &\text{ if $\sigma$ is odd}
        \end{cases}
    \]
    Thinking of $\{ 1,-1 \}$ as a group under multiplication, $\text{sgn}$ is a homomorphism. Since $1$ is the identity in $\{ 1,-1 \}$, $\text{ker(sgn)} = \text{sgn}^{-1} [\{ 1 \}]$ is a subgroup of $S_n$ that consists of all even permutations, which is exactly $A_n$. By the First Isomorphism Theorem, $S_n / A_n \cong \{ 1,-1 \}$, so the index of $A_n$ in $S_n$ is 2. Therefore, the order of $A_n$ is $n!/2$.
}

\Note{
    The set of all odd permutations has order $\displaystyle \frac{n!}{2}$ as well, but it does not form a subgroup since it is not closed under multiplication (the product of two odd permutations is an even permutation).
}

\Definition{Alternating Group}{
    The subgroup of $S_n$ consisting of all even permutations on $n$ letters is the \textbf{alternating group} $A_n$ on $n$ letters.
}

\Note{
    Both $S_n$ and $A_n$ are among the most important groups in mathematics, for two deep reasons.

    \textbf{Cayley's Theorem.} Every finite group $G$ is isomorphic to a subgroup of $S_n$ for
    $n = |G|$. The idea is that each element $g \in G$ induces a permutation of $G$ itself via
    left multiplication ($x \mapsto gx$), and collecting all such permutations yields a subgroup
    of $S_n$ isomorphic to $G$. In this sense, $S_n$ is the \emph{universal home} for all finite
    groups of order $\leq n$ -- permutation groups are not merely examples of groups, they are,
    up to isomorphism, \emph{all} groups.

    \textbf{Insolubility of the quintic.} While polynomial equations of degree $\leq 4$ admit
    closed-form solutions in terms of radicals, no such formula exists for degree $n \geq 5$.
    This is explained by Galois theory: a polynomial is solvable by radicals if and only if its
    Galois group is \emph{solvable} (i.e., built up from abelian pieces via a subnormal series).
    For $n \geq 5$, the alternating group $A_n$ is \textbf{simple} -- it admits no proper normal
    subgroups -- and is non-abelian, making it an indecomposable, unsolvable block. Since $A_n$
    embeds into the Galois group of a general degree-$n$ polynomial, the Galois group is not
    solvable, and no radical formula can exist.

    It is remarkable that a purely structural property of $A_n$ -- the absence of normal
    subgroups -- turns out to be the precise obstruction to solving polynomial equations by radicals.
}
